\documentclass[10pt,twocolumn]{article}

% ---- Packages (all in Overleaf's default TeX Live) ----
\usepackage[utf8]{inputenc}
\usepackage[T1]{fontenc}
\usepackage[margin=0.8in]{geometry}
\usepackage{booktabs}
\usepackage{amsmath,amssymb}
\usepackage{graphicx}
\usepackage{xcolor}
\usepackage{authblk}
\usepackage{titlesec}
\usepackage{tikz}
\usetikzlibrary{positioning,arrows.meta,fit,backgrounds,calc,shapes.geometric}
\usepackage[hidelinks]{hyperref}

\titleformat{\section}{\normalfont\large\bfseries}{\thesection}{0.5em}{}
\titleformat{\subsection}{\normalfont\normalsize\bfseries}{\thesubsection}{0.5em}{}
\titleformat{\subsubsection}{\normalfont\small\bfseries}{\thesubsubsection}{0.5em}{}
\setlength{\parskip}{2pt}

\newcommand{\rho}{\rho}
\newcommand{\code}[1]{\texttt{\small #1}}

\title{\textbf{Parallel Phase-Conditioned Cascading Memory Routing:\\
Two-Dimensional Concurrency for Multi-Agent LLM Memory\\
without Sacrificing Accuracy}}
\author[1]{[Your Name]}
\author[1]{[Mentor Name]}
\affil[1]{[Affiliation] \quad \texttt{[email]}}
\date{}

\begin{document}
\maketitle

\begin{abstract}
\noindent
The Phase-Conditioned Cascading Memory Router (PCCR) routes reads and writes
across six cognitive memory types over a ten-phase agentic lifecycle, gating
each optional store by a cost-effectiveness ratio
$\rho=\mathrm{utility}/\mathrm{cost}$. PCCR reduces \emph{retrieval cost}, but
its original execution model is strictly sequential: one task at a time, one
sub-agent at a time. This paper extends PCCR with \textbf{two orthogonal
dimensions of concurrency}---(1)~\emph{parallel sub-agents}, in which
independent sub-tasks within a single task run simultaneously, and
(2)~\emph{parallel tasks}, in which many tasks run simultaneously across the
system---while leaving the $\rho$-gate router \emph{byte-for-byte unchanged} and
preserving the ten-phase lifecycle. We achieve this with a seven-layer design
whose key invariants are per-task Working-Memory isolation (copy-on-fork), a
dependency-wave execution model with copy-on-read snapshots and a deterministic
staging$\to$merge, per-resource locks on the shared stores, and an exclusive
consolidation window. We prove---and empirically verify on 60 tasks
(\textbf{60/60 byte-identical} outcomes and identical final memory state)---that
the parallel schedule is \emph{outcome-equivalent} to the sequential one: it
changes \emph{latency only}, never accuracy. On a 100-task workload with a real
backbone latency profile (gpt-oss-120b, measured median $0.516$\,s/call),
parallel execution is \textbf{$3.73\times$} faster than sequential at
concurrency $C{=}4$ ($227.7\,\mathrm{s}\!\to\!61.1\,\mathrm{s}$, $-73\%$
wall-clock) and scales to \textbf{$13.5\times$} at $C{=}16$. The result is a
memory architecture that is simultaneously cost-efficient (via $\rho$-gating)
and throughput-efficient (via concurrency), with accuracy held constant.
\end{abstract}

% =====================================================================
\section{Introduction}
% =====================================================================
Memory-augmented LLM agents accumulate several functionally distinct memory
stores: procedural rules, volatile working state, a short-term plan cache,
episodic task traces, distilled semantic facts, and durable entity
profiles~\cite{coala,survey1}. PCCR~\cite{pccr} treats ``which store to consult''
as a \emph{decision} rather than a reflex: it conditions the policy on the
\emph{lifecycle phase} and the \emph{task pattern}, and on a cache miss gates
each optional store by an explicit ratio $\rho=U/C$ against a threshold
$\theta$. This cuts store consultations by 41--87\% at equal accuracy.

PCCR optimizes \emph{cost per task}. It says nothing about \emph{throughput} or
\emph{wall-clock latency}, because its reference execution model is sequential
in two senses at once: within a task the orchestrator delegates to one
sub-agent at a time, and across the system tasks are processed one after
another. Real office-automation workloads are neither: a single ``schedule a
sync and email everyone'' task contains several independent information-gathering
steps that need not be serialized, and a deployed system faces a stream of
tasks that could overlap. Serial execution leaves both forms of latency on the
table.

The central tension is that memory makes parallelism dangerous.
Working Memory is a per-task scratchpad that agents read and write throughout a
task; if two sub-agents mutate it concurrently they race. The shared
stores (STM/EM/SM/ENT) are written at task completion; if two tasks write
concurrently they can lose updates. And consolidation rewrites the FAISS
indices in bulk; if a task reads them mid-rewrite it sees a torn state. A naive
``just add threads'' approach would change task \emph{outcomes}, destroying the
very property---accuracy parity with retrieve-everything---that justifies PCCR.

\paragraph{Contributions.}
\begin{enumerate}\itemsep1pt
\item \textbf{A two-dimensional concurrency model} for phase-structured agent
memory: parallel sub-agents (intra-task) and parallel tasks (inter-task),
realized as seven cleanly separated layers (\S\ref{sec:arch},
Fig.~\ref{fig:arch}).
\item \textbf{Three correctness invariants}---per-task WM isolation, deterministic
dependency-wave execution, and per-resource locking with an exclusive
consolidation window---that together make the parallel schedule
\emph{outcome-equivalent} to the sequential one (\S\ref{sec:correct}).
\item \textbf{A non-invasive integration}: the $\rho$-gate router is unchanged,
the ten phases are preserved, the six stores gain only lock parameters, and all
prior tests still pass (\S\ref{sec:integration}).
\item \textbf{Empirical evidence} that parallelism buys latency without costing
accuracy: $3.73\times$ speedup at $C{=}4$ scaling to $13.5\times$ at $C{=}16$,
and \textbf{60/60 byte-identical} outcomes versus sequential
(\S\ref{sec:results}).
\end{enumerate}

% =====================================================================
\section{Background: the PCCR core}\label{sec:background}
% =====================================================================
We summarize only what is needed to follow the parallel extension; full detail
is in~\cite{pccr}.

\paragraph{Six memory types.} Table~\ref{tab:types} lists the stores and their
relative access cost. Procedural (PM), Working (WM), Short-term (STM), and
Entity (ENT) memory are cheap; Episodic (EM) and Semantic (SM) memory require a
vector (FAISS) search and are expensive. WM is special: it is \emph{per task}
and is unconditionally cleared when the task ends (``WM dies now'').

\begin{table}[t]
\centering\small
\begin{tabular}{@{}llc@{}}
\toprule
Type & Role & Access cost \\
\midrule
PM (proc.)    & static prompts/rules        & cheap (dict) \\
WM (work.)    & per-task scratchpad         & cheap; \emph{per task} \\
STM (short)   & cache of plan bundles       & cheap \\
EM (episodic) & past task traces (FAISS)    & \textbf{expensive} \\
SM (semantic) & distilled facts (FAISS)     & \textbf{expensive} \\
ENT (entity)  & per-user profiles           & cheap \\
\bottomrule
\end{tabular}
\caption{The six memory types. EXT (external files) is a seventh, non-memory
sink for artifacts (e.g.\ \code{.eml}/\code{.ics}).}
\label{tab:types}
\end{table}

\paragraph{Ten phases.} A task flows through:
(1)~bootstrap, (2)~ingestion, (3)~retrieval, (4)~orchestrator inference,
(5)~delegation, (6)~agent inference, (7)~agent output, (8)~observation,
(9)~storage, and (10)~consolidation. Phases 4--8 form the inner loop:
the orchestrator picks a sub-agent, the sub-agent acts, the environment
returns an observation, and the loop repeats until the task is done.

\paragraph{The $\rho$-gate (phase 3).} PM and WM are read unconditionally; STM
is checked first and a confident cache hit short-circuits all optional stores;
on a miss, an optional store $s\in\{\mathrm{EM},\mathrm{SM},\mathrm{ENT}\}$ is
consulted iff $\rho(\text{pattern},s)=U(\text{pattern},s)/C(s)\ge\theta$.
\textbf{Crucially, the entire routing decision happens once per task, in phase
3, before any sub-agent runs.} This single fact is what makes parallelism
orthogonal to routing: the execution model can be reorganized freely after
phase 3 without touching the router.

% =====================================================================
\section{The Parallel Architecture}\label{sec:arch}
% =====================================================================
We add concurrency as seven layers (Fig.~\ref{fig:arch}). Layers~1--2 introduce
parallel \emph{tasks}; layers~3--4 introduce parallel \emph{sub-agents};
layers~5--7 are the shared-state protections that make both safe.

% ----------- the big architecture figure -----------
\begin{figure*}[t]
\centering
\begin{tikzpicture}[
  font=\scriptsize,
  box/.style   ={draw, rounded corners, align=center, inner sep=3pt, minimum height=6.5mm},
  phase/.style ={draw, fill=gray!8,  rounded corners, align=center, inner sep=3pt, minimum height=6.5mm, text width=33mm},
  group/.style ={draw, fill=green!7, rounded corners, align=center, inner sep=4pt, text width=52mm},
  store/.style ={draw, fill=blue!10, rounded corners, align=center, minimum width=12mm, minimum height=5mm},
  storebox/.style={draw, fill=blue!4, rounded corners, align=center, inner sep=5pt, text width=33mm},
  cons/.style  ={draw, fill=orange!12, rounded corners, align=center, inner sep=4pt, text width=70mm},
  lyr/.style   ={font=\scriptsize\bfseries, text=red!55!black},
  arr/.style   ={-{Latex[length=1.6mm]}, semithick},
  darr/.style  ={-{Latex[length=1.6mm]}, dashed, semithick},
  node distance=3.2mm and 9mm
]

% ---- LAYER 1: task queue ----
\node[box, fill=orange!12, text width=150mm] (queue)
  {\textbf{LAYER 1 — Async Task Queue}\quad accepts $N$ tasks, runs their pipelines concurrently
   (\code{max\_concurrency}=$C$); admits a task into phase 3 only when \code{is\_consolidating()} is false};

% ---- LAYER 2: N per-task pipelines (show 3) ----
\node[box, below left=5mm and 18mm of queue, text width=24mm] (t1) {Task 1 pipeline\\\textit{own WM}};
\node[box, below=5mm of queue, text width=24mm]               (t2) {Task 2 pipeline\\\textit{own WM}};
\node[box, below right=5mm and 18mm of queue, text width=24mm](tn) {Task $N$ pipeline\\\textit{own WM}};
\node[lyr, left=1mm of t1] {L2};
\draw[arr] (queue) -- (t1); \draw[arr] (queue) -- (t2); \draw[arr] (queue) -- (tn);

% ---- expanded single pipeline (center column) ----
\node[phase, below=8mm of t2] (p12) {\textbf{Phase 1--2}\\ bootstrap (shared, once) + ingestion\\ task$\to$WM, pattern classified};
\node[phase, below=of p12]    (p3)  {\textbf{Phase 3 — PCCR $\rho$-gate}\\ PM+WM always; STM check;\\ $\rho=U/C\ge\theta\Rightarrow$ consult $s$ \emph{(once/task)}};
\node[phase, below=of p3]     (p45) {\textbf{Phase 4--5} orchestrator\\ emit next \emph{group} of delegations};
\node[box, fill=gray!8, below=of p45, text width=33mm] (dag)
  {\textbf{L3 — Dependency Analyzer}\\ group $\to$ DAG; indep.\ $\to$ same wave,\\ dependent $\to$ next wave};
\node[group, below=of dag] (grp)
  {\textbf{L4 — Parallel agent wave}\\[1pt]
   1.\ snapshot WM (frozen, copy-on-read)\\
   2.\ \fbox{calendar agent} $\parallel$ \fbox{email agent} $\parallel\,\dots$\\
   \hspace*{3mm}(each writes its own \emph{staging} slot)\\
   3.\ \code{asyncio.gather} runs the wave\\
   4.\ deterministic merge (sorted) $\to$ live WM};
\node[phase, below=of grp] (p8) {\textbf{Phase 8} merge staging $\to$ WM\\ \code{step\_history}, \code{shared\_context}};
\node[draw, diamond, aspect=2.2, fill=yellow!15, below=of p8, inner sep=1pt, align=center] (dec) {more\\groups?};
\node[phase, below=4mm of dec] (p9)
  {\textbf{Phase 9} — write routing (W1--W6)\\ success gate; pattern$\to$STM (lock);\\ person$\to$ENT (lock); EM/SM deferred;\\ routing history (atomic, W6); \textbf{WM clear}};

\draw[arr] (p12) -- (p3); \draw[arr] (p3) -- (p45); \draw[arr] (p45) -- (dag);
\draw[arr] (dag) -- (grp); \draw[arr] (grp) -- (p8); \draw[arr] (p8) -- (dec);
\draw[arr] (dec) -- node[right]{done} (p9);
\draw[arr] (dec.west) -- ++(-26mm,0) node[midway,above]{yes (re-plan)} |- (p45.west);

% pipeline container (background)
\begin{scope}[on background layer]
\node[draw, dashed, rounded corners, fill=black!2, fit=(p12)(p3)(p45)(dag)(grp)(p8)(dec)(p9),
      inner sep=4mm, label={[lyr]above right:{Per-task pipeline (Task $i$) — own 10-phase lifecycle, own WM \;$\times N$ concurrent}}] (pipe) {};
\end{scope}

% ---- LAYER 5/6: shared stores (right) ----
\node[storebox, right=14mm of p3, anchor=west] (stores)
  {\textbf{LAYER 5 — Shared stores}\\ \textbf{(cross-task, lock-protected)}\\[2pt]
   PM \;\; STM\textsuperscript{$\dagger$} \;\; ENT\textsuperscript{$\dagger$}\\
   EM \;\; SM \;\; routing history\\[2pt]
   \footnotesize $\dagger$ per-key lock (per pattern / per user);\\
   EM/SM/PM via global FAISS/PM lock};
\node[lyr, above=0.5mm of stores] {L5/L6};
\draw[arr] (p3.east) to[bend left=8] node[above,sloped]{read} (stores.west);
\draw[arr] (p9.east) to[bend right=14] node[below,sloped]{write (locked)} (stores.south west);

% ---- LAYER 7: consolidation (bottom) ----
\node[cons, below=8mm of p9] (p10)
  {\textbf{LAYER 7 — Phase 10 Consolidation (exclusive window)}\\
   raises \code{is\_consolidating}; holds global FAISS lock; batch-writes EM/SM/PM/STM.\\
   \emph{No task enters phase 3 while open.}};
\draw[arr] (p9) -- (p10);
\draw[arr] (p10.east) to[bend right=18] node[below,sloped]{rewrites} (stores.south);
\draw[darr] (p10.west) -- ++(-20mm,0) |- node[near start,left]{blocks}([xshift=-6mm]queue.west);

\end{tikzpicture}
\caption{The seven-layer parallel PCCR architecture. \textbf{Parallel tasks}
(L1--L2): the async queue runs $N$ task pipelines concurrently, each with its
\emph{own} Working Memory (copy-on-fork) but sharing the cross-task stores.
\textbf{Parallel sub-agents} (L3--L4): within a task the orchestrator emits a
group of delegations; the dependency analyzer splits it into waves (independent
sub-tasks in one wave); a wave runs concurrently against a frozen WM snapshot
and is merged back deterministically; control returns to the orchestrator
(``more groups?'') until \textsc{finish}. \textbf{Protections} (L5--L7):
per-resource locks guard concurrent store writes, and phase 10 is an exclusive
window. The $\rho$-gate (phase 3) is unchanged and runs once per task before any
parallelism.}
\label{fig:arch}
\end{figure*}

\subsection{Design goals and invariants}
The extension is governed by four hard constraints:
(C1)~the $\rho$-gate router is unmodified; (C2)~the ten-phase lifecycle is
preserved; (C3)~the six stores change only by gaining optional lock parameters;
(C4)~the parallel schedule yields the \emph{same task outcomes} as the
sequential one. C1--C3 keep the contribution surgical; C4 is the correctness
property proved in \S\ref{sec:correct}. Three invariants deliver C4:

\begin{description}\itemsep1pt
\item[\textsc{I1 (WM isolation).}] No two concurrently executing units (tasks
or same-wave agents) ever mutate the same Working Memory.
\item[\textsc{I2 (deterministic merge).}] The result of merging a wave's
sub-agent outputs into WM is independent of the order in which the sub-agents
finish.
\item[\textsc{I3 (no lost writes).}] Concurrent writes to a shared store are
serialized per resource, and consolidation never overlaps task reads.
\end{description}

\subsection{Layer 1 --- Async task queue (parallel tasks)}
\code{AsyncTaskQueue} accepts a list of task specs and runs each task's full
pipeline as an \code{asyncio} coroutine, bounded by a semaphore of size $C$
(\code{max\_concurrency}). Two behaviours matter. First, \emph{fault isolation}:
a task that raises is caught and returns an error record, so one failure cannot
abort the batch. Second, the \emph{consolidation gate}: before a task enters
phase 3 (the FAISS read), the queue waits while \code{is\_consolidating()} is
true, so no task reads vectors that phase 10 is rewriting (this realizes half of
\textsc{I3}). The queue reduces \emph{inter-task} latency: while task~$A$ waits
on an LLM call, task~$B$ proceeds.

\subsection{Layer 2 --- Per-task pipeline with WM isolation}
Each task runs its own ten-phase lifecycle. The mechanism is
\code{fork\_for\_task()}: it produces a per-task \code{MemoryManager} that
\emph{shares} the cross-task stores (PM/STM/EM/SM/ENT), the router, and the lock
manager, but receives a \emph{fresh} Working Memory instance. This is the
copy-on-fork that satisfies \textsc{I1} at the task level: because WM is per task
by construction (Table~\ref{tab:types}), giving each concurrent task its own WM
means concurrent tasks can never collide on working state. The bootstrap (phase
1) is done once on the parent and shared, not repeated per task.

\subsection{Layer 3 --- Dependency analyzer (DAG $\to$ waves)}
Within a task, the orchestrator emits a \emph{group} of delegations
(an agent + a sub-task instruction each). The analyzer turns a flat group into
ordered \emph{waves}. It assigns each delegation a coarse category by keyword---
\textsc{read} (list/search/find/\ldots), \textsc{create} (create/add/schedule/
\ldots), \textsc{send} (send/email/notify/\ldots), or \textsc{finish}---and
orders categories by an action-precedence relation
$\textsc{read}\prec\textsc{create}\prec\textsc{send}\prec\textsc{finish}$.
Delegations in the same category form one wave (run in parallel); categories run
in sequence. The intuition is sound and conservative: reads gather information
that creates may depend on, and sends (notifications) depend on what was
created. Formally, every node in wave $k{+}1$ is treated as depending on every
node in wave $k$ (a coarse but \emph{safe} over-approximation of the true DAG:
it never parallelizes a genuine dependency). The motivating example
``list events $\to$ create event $\to$ send invite'' becomes three ordered
waves; three independent reads become a single three-way parallel wave.

\subsection{Layer 4 --- Parallel agent execution}
A wave executes in four steps that together guarantee \textsc{I1} (agent level)
and \textsc{I2}:
\begin{enumerate}\itemsep0pt
\item \textbf{Snapshot.} \code{wm.snapshot()} deep-copies the current WM. Every
agent in the wave reads this \emph{frozen} copy, never live WM (copy-on-read).
\item \textbf{Staging.} Each agent writes only its own result object, never live
WM.
\item \textbf{Gather.} \code{asyncio.gather} runs the wave concurrently; the
overlap of the agents' (I/O-bound) LLM calls is where wall-clock is saved.
\item \textbf{Merge.} After \emph{all} agents finish, results are merged into
live WM in a fixed sorted order (by \code{(agent\_type, sub\-task)}), appending
one step per agent and updating \code{shared\_context}.
\end{enumerate}
Because same-wave agents are \emph{independent} (L3 placed them together
precisely because neither depends on the other), they read external state, not
each other; the frozen snapshot is therefore sufficient, and the sorted merge
makes the outcome invariant to completion order (\textsc{I2}). Dependent agents
are in different waves, so a later agent does see the merged results of earlier
waves.

\subsection{Layer 5 --- Memory lock manager}
\code{MemoryLockManager} hands out \code{asyncio} locks at a granularity matched
to contention. STM is locked \emph{per pattern key} and ENT \emph{per username},
so two tasks touching different patterns/users never block each other; only
genuine same-key contention serializes. EM/SM share a single global FAISS lock
and PM a single PM lock, because their writes are global, batched, and rare
(phase 10). A \code{history\_lock} guards the shared router's routing-history
append (W6). Locks are created lazily on first use, because on the target
runtime an \code{asyncio.Lock} must bind to a running event loop. This layer
delivers \textsc{I3}.

\subsection{Layer 6 --- Phase-9 parallel writes}
Phase 9 (storage) applies six write rules; the parallel version runs the
success-gated writes concurrently via \code{asyncio.gather}, each acquiring its
own lock:
W1~success gate (failures write nothing);
W2~pattern$\to$STM under the per-pattern lock;
W3~person$\to$ENT under the per-user lock;
W4~experience$\to$EM and W5~rule$\to$PM are \emph{deferred} to consolidation
(only the raw trace is appended now);
W6~routing-history append under \code{history\_lock} (atomic across the tasks
that share one router). WM is then cleared. Because the writes touch disjoint,
independently locked resources, they proceed in parallel without lost updates.

\subsection{Layer 7 --- Consolidation exclusive window}
Consolidation (phase 10) distills the accumulated traces into EM/SM, learns PM
rules, and pre-fills STM---a bulk rewrite of the FAISS indices. The async
wrapper \code{consolidate\_async} raises \code{is\_consolidating()} and holds the
global FAISS lock for the duration, then lowers the flag in a \code{finally}
block. Combined with the L1 gate, this makes phase 10 an \emph{exclusive
window}: in-flight tasks finish, no new task enters phase~3, the indices are
rewritten atomically, and normal concurrency resumes. This is the second half of
\textsc{I3}.

\subsection{The orchestrator re-planning loop}
Layers 3--4 run \emph{one group}; a full task may need several. The loop
\code{run\_planned\_task} re-invokes the orchestrator after each group merges:
the orchestrator inspects the (now-updated) live WM, emits the next group (or
\textsc{finish}), that group is split into waves and run, and control returns
(the ``more groups?'' decision in Fig.~\ref{fig:arch}). This keeps the
orchestrator in the loop---each round sees the merged results of all prior
groups---rather than decomposing the whole plan up front.

% =====================================================================
\section{Correctness: parallel $\equiv$ sequential}\label{sec:correct}
% =====================================================================
We argue that the parallel schedule produces the same outcome as the sequential
one, then verify it empirically (\S\ref{sec:eq}).

\paragraph{Intra-task (parallel sub-agents).}
Consider any task. Sequential execution runs its delegations one at a time in
dependency order; parallel execution runs each wave concurrently. The two differ
\emph{only} in whether same-wave agents observe one another's mid-flight WM
writes. By L3, same-wave agents are mutually independent, so an agent's result
is a function of external state and its own sub-task, \emph{not} of a sibling's
WM step. Hence each agent produces the same result in both schedules. The merge
(L4) is order-independent (\textsc{I2}). Dependent agents are in different waves,
and a later wave reads the merged earlier waves identically in both schedules.
Therefore the merged WM---and thus the task's plan, actions, and
success---are identical.

\paragraph{Inter-task (parallel tasks).}
Each task has its own WM (\textsc{I1}), so concurrent tasks cannot corrupt one
another's working state. Shared-store \emph{writes} are serialized per resource
(\textsc{I3}); no update is lost. The one schedule-dependent quantity is what a
task \emph{reads} from the shared cache: under concurrency, task~$B$ may run
before task~$A$'s STM write lands and therefore \emph{miss} a cache entry it
would have hit sequentially. But a miss only causes a \emph{recompute}, and the
recomputed plan equals the cached plan (the cache stores a previously computed
plan for the same signature). Thus cross-task concurrency changes cache-hit
\emph{timing}---an efficiency quantity (consultations, latency)---never the
correctness of any task's answer. Standard held-out evaluation scores each task
independently, so per-task accuracy is unaffected.

\paragraph{The one genuine divergence vector.}
The argument assumes each task acts on its \emph{own} external environment. If
tasks share a single mutable external resource (e.g.\ one global filesystem of
calendar events), a concurrent task can mutate it under another task's read,
producing a different observation. This is a property of the shared
\emph{environment}, not of the memory architecture, and held-out evaluation
already isolates it (each task is scored against its own sandbox)---the external
analogue of per-task WM isolation. \S\ref{sec:eq} demonstrates both the
divergence (under a deliberately shared world) and its disappearance under
per-task isolation.

\paragraph{Unit-tested primitives.} The invariants are pinned by tests:
deterministic merge regardless of finish order (\textsc{I2}); snapshot isolation
(\textsc{I1}); per-key lock identity and no-lost-update under contention, and the
consolidation-window wait (\textsc{I3}); plus fault isolation and the
re-planning loop. All 39 tests pass (15 exercise the parallelism layer).

% =====================================================================
\section{Integration (non-invasive)}\label{sec:integration}
% =====================================================================
The parallel path is \emph{additive}. New modules \code{parallel.py} (analyzer,
executor, lock manager) and \code{task\_queue.py} are self-contained. The
manager gains async methods (\code{run\_planned\_task},
\code{complete\_task\_async}, \code{consolidate\_async}, \code{fork\_for\_task})
that sit \emph{beside} the unchanged synchronous lifecycle methods; the stores
gain only \code{snapshot}/\code{merge\_staging} on WM. The $\rho$-gate router
file is untouched (C1), the ten phases keep their identity (C2), and every
pre-existing test still passes (C3/C4). A system can therefore run sequentially
or in parallel from the same codebase.

% =====================================================================
\section{Experiments}\label{sec:results}
% =====================================================================
\subsection{Setup}
The benchmark drives the \emph{same} real framework (per-task fork, $\rho$-gate,
re-planning loop, dependency waves, locked writes, shared stores) in two
execution modes---sequential and parallel---over a 100-task workload with a
realistic pattern mix (single-action 20, lookup 20, recurring 10, coordination
30, exploratory 20; 440 LLM calls in total, critical-path depth 360). The only
modeled quantity is the external call cost $L$, set to the \textbf{measured}
per-call latency of gpt-oss-120b on Cerebras: median $0.516$\,s over real calls
(mean $0.743$\,s). Because \code{asyncio} overlaps a real network await exactly
as it overlaps \code{asyncio.sleep($L$)}, the measured wall-clock speedup is what
real parallel API calls would achieve for this I/O-bound workload.

\subsection{Latency: parallel vs.\ sequential}
At concurrency $C{=}4$, parallel execution completes the 100 tasks in
$61.1$\,s versus $227.7$\,s sequential---a \textbf{$3.73\times$ speedup}, $-73\%$
wall-clock (Table~\ref{tab:lat}). The sequential time matches the analytic
$440{\times}0.516{=}227.0$\,s, confirming the model.

\begin{table}[t]
\centering\small
\begin{tabular}{@{}lccc@{}}
\toprule
Run & Wall-clock (100 tasks) & Per task & Speedup\\
\midrule
Sequential       & $227.7$\,s & $2277$\,ms & $1.00\times$\\
\textbf{Parallel ($C{=}4$)} & $\mathbf{61.1}$\,\textbf{s} & $\mathbf{611}$\,\textbf{ms} & $\mathbf{3.73\times}$\\
\bottomrule
\end{tabular}
\caption{Latency at the measured backbone latency ($L{=}0.516$\,s/call).}
\label{tab:lat}
\end{table}

\subsection{Scaling with concurrency}
The speedup grows near-linearly with $C$ up to the dependency critical path
(Table~\ref{tab:scale}). The ratio is latency-independent: $C{=}4$ measured at a
reduced $L{=}0.05$\,s reproduces the same $3.7\times$, confirming the framework
overhead is negligible relative to call latency. $C{=}1$ is approximately neutral
($0.96\times$): with no task-level concurrency the only gain is intra-task wave
overlap, which roughly cancels the per-wave snapshot/merge overhead. The win is
therefore dominated by \emph{parallel tasks}, with \emph{parallel sub-agents} a
secondary benefit on coordination/exploratory tasks.

\begin{table}[t]
\centering\small
\begin{tabular}{@{}lccccc@{}}
\toprule
$C$ & 1 & 2 & 4 & 8 & 16\\
\midrule
Speedup        & $0.96\times$ & $1.89\times$ & $3.73\times$ & $7.15\times$ & $13.47\times$\\
Wall-clock cut & $-5\%$ & $47\%$ & $73\%$ & $86\%$ & $93\%$\\
\bottomrule
\end{tabular}
\caption{Speedup vs.\ task-concurrency $C$ (100 tasks).}
\label{tab:scale}
\end{table}

\subsection{Accuracy equivalence}\label{sec:eq}
We run the same task set through both schedules and compare, per task, the
outcome tuple (success, plan, sorted normalized observations) and the final
shared-memory state (STM/EM/SM/PM sizes and ENT facts). With \textbf{per-task
isolated environments}, outcomes are \textbf{60/60 byte-identical} and the final
memory state is identical (Table~\ref{tab:eq})---empirically confirming
\S\ref{sec:correct}. With a deliberately \emph{shared} external world, 2/60
coordination tasks diverge on an incidental observation count (the number of
calendar events seen by a search agent), while success, plan, and memory state
remain identical---exactly the single divergence vector predicted above, and
removed by isolation.

\begin{table}[t]
\centering\small
\begin{tabular}{@{}lcc@{}}
\toprule
Environment & Outcomes identical & Memory state identical\\
\midrule
Per-task isolated & \textbf{60/60} & \textbf{yes}\\
Shared world      & 58/60\textsuperscript{$\ast$} & yes\\
\bottomrule
\end{tabular}
\caption{Parallel vs.\ sequential outcome equivalence ($C{=}4$, 60 tasks).
\textsuperscript{$\ast$}The two differ only on an incidental observation count
of a shared mutable resource; success/plan/memory identical.}
\label{tab:eq}
\end{table}

\subsection{Cost/accuracy carry over unchanged}
Because the $\rho$-gate is untouched and outcomes are equivalent, the prior
PCCR cost and accuracy results~\cite{pccr}---41--87\% fewer consultations and up
to 89\% fewer injected tokens at equal accuracy on OfficeBench---hold verbatim
for the parallel system. Parallelism and $\rho$-gating optimize orthogonal axes:
the former cuts \emph{latency}, the latter cuts \emph{cost per task}; neither
moves accuracy.

% =====================================================================
\section{Related Work}
% =====================================================================
\paragraph{Parallel multi-agent execution.} DynTaskMAS~\cite{dyntaskmas}
introduces a shared context layer (SACMS) and asynchronous parallel execution;
GAP~\cite{gap} models sub-task dependencies as a DAG; SPOQ~\cite{spoq} dispatches
in topological waves; AdaptOrch~\cite{adaptorch} switches between parallel and
sequential executors. We adopt the DAG-into-waves idea (L3) and async dispatch
(L1, L4) but place them \emph{underneath a phase-conditioned memory router},
which prior parallel orchestrators do not have. CompArch~\cite{comparch}
motivates copy-on-read snapshots for memory consistency (L4);
AISAC~\cite{aisac} uses a blackboard append/patch merge (our staging$\to$merge).

\paragraph{Concurrent memory systems.} MIRIX~\cite{mirix} manages six memory
types with \emph{separate per-type agents} and uses \code{asyncio.gather} with
per-resource locks for parallel writes; we borrow the locking pattern (L5--L6)
but keep a \emph{single} central $\rho$-gated manager rather than per-type
agents, and we add the inter-task queue and the exclusive consolidation window.
LEGOMem~\cite{legomem} splits procedural memory across orchestrator/sub-agent
modules in a multi-agent setting but is sequential and has no cost gate.

\paragraph{Cost-sensitive memory routing.} The PCCR line~\cite{pccr,pocket}
treats retrieval as a cost-sensitive decision; none of these address execution
concurrency. To our knowledge, this is the first system to combine a six-type
cost-gated memory router with a two-dimensional concurrency model and a proof of
outcome-equivalence.

% =====================================================================
\section{Limitations}
% =====================================================================
\textbf{Latency model.} We model agent work as the measured per-call latency
rather than re-issuing 440 live calls (cost/rate-limit); the \emph{scheduling}
measured is the real framework's, but absolute numbers assume the I/O-bound
regime. \textbf{Provider limits.} Effective $C$ is capped below the ideal by
provider rate limits (we observed HTTP~429 at four keys) and by the dependency
critical path (a \textsc{send} still waits for its \textsc{create}).
\textbf{Dependency granularity.} L3 uses a coarse, keyword-based,
category-precedence DAG; it is safe (never parallelizes a true dependency) but
may serialize some independent work (false dependencies), leaving speedup on the
table. A learned or semantic dependency analyzer is future work.
\textbf{Environment isolation.} Outcome-equivalence assumes per-task external
isolation; systems that share a live mutable environment must add
application-level resource coordination, exactly as benchmarks already do.

% =====================================================================
\section{Reproducibility}
% =====================================================================
The parallelism lives in \code{memory\_manager/parallel.py} and
\code{task\_queue.py} with manager hooks in \code{manager.py}. Three scripts
reproduce the paper: \code{run\_parallel\_demo.py} (all seven layers end-to-end,
offline), \code{bench\_parallel.py} (Tables~\ref{tab:lat}--\ref{tab:scale}), and
\code{verify\_equivalence.py} (Table~\ref{tab:eq}). All 39 unit tests run with
\code{pytest}. No API key is needed for the demo, benchmark, or equivalence
check (deterministic stub backends); the latency constant is the only measured
input and is overridable via \code{-{}-latency}.

% =====================================================================
\section{Conclusion}
% =====================================================================
We extended PCCR with two orthogonal dimensions of concurrency---parallel
sub-agents within a task and parallel tasks across the system---through a
seven-layer design that leaves the $\rho$-gate router and the ten-phase
lifecycle intact. Per-task Working-Memory isolation, deterministic
dependency-wave execution, and per-resource locking with an exclusive
consolidation window make the parallel schedule provably and empirically
outcome-equivalent to the sequential one: $3.73\times$ faster at $C{=}4$
(up to $13.5\times$ at $C{=}16$) with \textbf{60/60 byte-identical} outcomes.
A memory architecture can be made cost-efficient and throughput-efficient at the
same time, with accuracy held exactly constant.

\begin{thebibliography}{99}\small
\bibitem{pccr} [Author(s)]. Phase-Conditioned Cascading Memory Routing for Multi-Agent LLM Systems. Companion manuscript, 2026.
\bibitem{coala} T.~Sumers et al. Cognitive Architectures for Language Agents (CoALA). TMLR, 2024.
\bibitem{survey1} Memory for Autonomous LLM Agents: Mechanisms, Evaluation, and Emerging Frontiers. arXiv:2603.07670.
\bibitem{dyntaskmas} DynTaskMAS: A Dynamic Task-Graph Multi-Agent System with Asynchronous Parallel Execution. arXiv:2503.07675.
\bibitem{gap} GAP: Graph-based Agent Planning with Dependency-Aware Parallelism. arXiv:2510.25320.
\bibitem{spoq} SPOQ: Scheduling Parallel Orchestration via Topological Wave Dispatch. arXiv:2606.03115.
\bibitem{comparch} A Compositional Architecture for Consistent Agent Memory. arXiv:2603.10062.
\bibitem{aisac} AISAC: A Blackboard-Based Shared-Context Substrate for Multi-Agent Systems. arXiv:2511.14043.
\bibitem{adaptorch} AdaptOrch: Adaptive Parallel/Sequential Orchestration for LLM Agents. arXiv:2602.16873.
\bibitem{mirix} Y.~Wang, X.~Chen. MIRIX: Multi-Agent Memory System for LLM-Based Agents. arXiv:2507.07957.
\bibitem{legomem} LEGOMem: Modular Procedural Memory for Multi-agent LLM Systems. arXiv:2510.04851.
\bibitem{pocket} M.~Gaikwad. Did You Check the Right Pocket? Cost-Sensitive Store Routing for Memory-Augmented Agents. arXiv:2603.15658.
\bibitem{officebench} T.~Wang et al. OfficeBench: Benchmarking Language Agents across Multiple Applications for Office Automation. arXiv:2407.19056.
\end{thebibliography}

\end{document}
